\paragraph{Identified model energies and the physical particle boundary.}
The model energies and inertial mass are identifiable.  The cube has eight
exactly characterized interacting energies, each reached by the local
probe with positive weight.  The moving pair has
$\epsilon(P)=|P|^2/4-\gamma^2$ for the reflected generator
$J_a=6/a^2-H_a$, hence binding energy $\gamma^2$ and
inertial mass~2 in those model units.  If physical energy and length units
$E_{\rm u},L_{\rm u}$ are independently attached and
$p_{\rm phys}=\hbar P/L_{\rm u}$, then
$M=2\hbar^2/(E_{\rm u}L_{\rm u}^2)$ and
$E_{\rm bind}=E_{\rm u}\gamma^2$, with
$E_{\rm u}T_{\rm u}=\hbar$ for the common clock.
These conversion formulas do not select the units or fix an absolute
rest-energy offset.  Binding energy is not total rest mass.
For the cube, the following rounded displays have exact rational
energy and weight intervals in the receipt, with the original contact
unit $U=1$:
\begingroup
\setlength{\tabcolsep}{8pt}
\begin{center}
\begin{tabular}{rrr}
$H$ energy & $E-3$ & Probe weight\\
\hline
1.816498920 & $-1.183501080$ & 0.002465366\\
1.860537501 & $-1.139462499$ & 0.005644624\\
1.955137632 & $-1.044862368$ & 0.005881219\\
2.108692846 & $-0.891307154$ & 0.007137664\\
2.167639734 & $-0.832360266$ & 0.004528123\\
3.015861346 & $+0.015861346$ & 0.421577939\\
3.030769653 & $+0.030769653$ & 0.415789141\\
3.044862368 & $+0.044862368$ & 0.136975924
\end{tabular}
\end{center}
\endgroup
Weights refer to the normalized centered density probe; the phase-kicked
target has $15/16$ times these weights plus an elastic weight $1/16$ at
energy~3.  The lowest displayed line is not the $N=6$ ground:
$-3\sqrt3/8\le E_{0,N=6}\le-\sqrt3/8$, by the free one-particle
bound and a polarized zero-contact variational state.  These are finite
energy levels and signed transitions, not eight rest masses.
The physical species assignment count remains zero.  On the cube the density
probe, its phase-kicked reference and all spectral components remain in
$N=6$: the transition operator is number-neutral and even under
$\Gamma=(-1)^N$.  An isometric dictionary intertwining this grading cannot
map that response to a single odd fermion, since an odd creation operator
changes parity.  This does not exclude other odd model excitations or a
fermion--antifermion response; even parity alone does not establish Bose
exchange or Lorentz spin.  Hidden spin and model number have not been
identified with physical spin and electric charge.
The physical target conventions are the PDG electron and quark quantum
numbers and relativistic kinematics.\footnote{\url{https://pdg.lbl.gov/2025/listings/rpp2025-list-electron.pdf};
\url{https://pdg.lbl.gov/2025/reviews/rpp2025-rev-quark-model.pdf};
\url{https://pdg.lbl.gov/2025/reviews/rpp2025-rev-kinematics.pdf}.}

The supplied cube reference also fails the vacuum ground-state spectral
condition for its own $H$.  Its positive density-probe spectral measure has
five negative and three positive frequencies $\omega=z/16-1$; all eight
weights are strictly positive by the exact Krylov certificate.  In
particular $z_-=8-2\sqrt{19}$ has weight $(1-4/\sqrt{19})/14>0$ and gives
an $N=6$ eigenenergy $5/2-\sqrt{19}/8<2$, below the reference energy~3.
A number-preserving unitary swapping these two states lowers the energy,
so the reference is not passive for the full field unitary algebra.
This is a mathematical witness, not an executed extraction protocol:
the actual local phase preserves the mean energy.  These are frequencies
of a positive greater correlation, not the two-sided frequencies of a
retarded commutator.  A scalar or common number chemical potential leaves
the gaps unchanged, and time reversal leaves mixed signs.

The moving-pair limit $\epsilon(P)=|P|^2/4-\gamma^2$ is a genuine
controlled nonrelativistic band.  For any fixed affine energy map
$E=\alpha\epsilon+\beta$, $\alpha\ne0$, and fixed invertible linear
momentum map $p=LP$, the residual $E^2-c^2|p|^2-m^2c^4$ has nonzero
quartic coefficient $\alpha^2/16$ along a coordinate line.  It cannot
vanish on an open momentum set.  Thus these maps do not turn this limit
into an exact massive or massless Lorentz shell.  A low-momentum
approximation remains possible: after matching constant and quadratic
terms, the difference between $E_0+Ar^2$ and $\sqrt{E_0^2+2E_0Ar^2}$ is
$A^2r^4/(E_0+Ar^2+\sqrt{E_0^2+2E_0Ar^2})$, for rest energy
$E_0>0$ and $A>0$.

\texttt{model/particle\_identification} binds these scoped tests and the
remaining vacuum, gauge-charge, spin, clock, detector and mass-shell
requirements.  The earlier dimer electron mass was a calibration input,
not an independent mass prediction; its retrospective equality comparison
does not identify the new cube or pair modes.  No mass fitting, new eligible
physical comparison, or whole-program exclusion is added.
